\chapter*{Preface}

Two gears have $42$ and $30$ teeth. A chalk mark on each comes back to the contact point after $210$ teeth have passed: the least positive multiple shared by $42$ and $30$. Euclid's algorithm finds the number behind this meeting: $\gcd(42,30)=6$. It takes two divisions, not a search through every possible meeting time.

That calculation is a model for this book. A question begins with numbers you can touch. A pattern suggests a claim. A proof explains why the claim survives when the numbers change. Sometimes the same method finds a solution. Sometimes it proves that none can exist.

Primes make the contrast plain. Every integer greater than $1$ is a product of primes in one way, apart from the order of the factors. Yet no finite list contains all the primes. We can prove that primes never run out and that one lies between $n$ and $2n$ for every $n>1$. Neither theorem tells us where the next prime will be. Later volumes will ask how accurately primes can be counted.

The arithmetic of remainders changes the question. The congruence $x^2\equiv-1\pmod p$ may have a solution for one odd prime and none for another. Quadratic reciprocity connects apparently separate decisions. Continued fractions turn rough approximations into good ones and find solutions of $x^2-Dy^2=1$. Gaussian integers turn $x^2+y^2=p$ into factorisation in a larger ring. Each change of language should solve a problem the old language left awkward.

\section*{What you need}

Volume~1, \emph{First Principles}, supplies the habits: read quantifiers in order, prove both directions of an equivalence, and check boundary cases. We use its induction, counting, polynomials, complex numbers, inequalities, limits, elementary calculus and linear algebra as needed. The first chapters require much less: integer arithmetic and a willingness to follow a proof.

No earlier course in abstract algebra is assumed. Groups, rings and ideals are defined when they first do a job. Volume~3 develops their general theory. Volume~13 treats general algebraic integers, ideals and $p$-adic fields. Here we prove unique factorisation for Gaussian and Eisenstein integers and show why it can fail in $\Z[\sqrt{-5}]$; we do not quietly assume the later remedy.

Advanced theorems mentioned for perspective are not hidden premises. The prime number theorem, general quadratic fields and the full $p$-adic construction are later destinations. The book distinguishes a statement proved here from an outlook.

\section*{Two routes through the book}

For a first reading, follow the chapter order. Begin with divisibility and primes. Continue through congruences, primitive roots, reciprocity and arithmetic functions. The final part treats sums of squares, approximation, Diophantine equations, complex integer rings and characters. Work through each example. Try the one-star exercises before their hints; leave a long challenge marked rather than letting it block the next chapter.

For a deeper reading, return to the two- and three-star problems. Some ask you to connect different chapters or use a theorem from another volume. Read that tool first, or treat the problem as a preview. The stars measure the work asked for, not the value of the mathematics.

When stuck, ask what the statement demands. For divisibility, exhibit a factor. For a congruence, name the modulus and check cancellation. For a proposed square, test small residues. For an equation, finding one solution is not the same as classifying all solutions. Compare a first example with an awkward one: zero, one, two, a negative input or a repeated prime factor.

The diagrams show choices in the proof: a remainder shrinking, a residue class surviving, or a fraction approaching a square root. They are evidence for a conjecture and pictures of an argument; they are not substitutes for the argument. Historical photographs are credited where they appear.

The \emph{Arithmetic Field Guide} at the end collects definitions and practical first moves. Dirichlet characters and Euler products close this volume by turning arithmetic conditions into sums. Volume~6 uses that language to count primes. Volume~13 asks what becomes of factorisation in larger rings. The integers stay the same; our ways of seeing them grow.
